\section{Suffix Array y LCP}
\whenuse{Ordenar todos los sufijos para comparar substrings, buscar patrones,
contar substrings distintos o conocer prefijos comunes sin usar hashing.}

\texttt{sa[i]} es la posición donde comienza el $i$-ésimo sufijo en orden
lexicográfico. \texttt{lcp[i]} es la longitud del prefijo común entre los
sufijos \texttt{sa[i]} y \texttt{sa[i+1]}. Esta versión usa doubling y
Con sa podés:
- Ordenar los sufijos.
- Buscar un patrón con búsqueda binaria.
- Comparar substrings.
- Encontrar repeticiones.
Con lcp podés:
- Encontrar el substring repetido más largo.
- Contar substrings distintos.
- Consultar cuánto coinciden dos sufijos.
\texttt{sort}.

\complexity{Suffix Array $O(n\log^2 n)$; LCP $O(n)$; memoria $O(n)$.}

\begin{lstlisting}[style=cpp]
vector<int> suffix_array(const string &s) {
    int n=s.size();
    vector<int> order(n),rank(n),new_rank(n);
    iota(order.begin(),order.end(),0);
    for(int i=0;i<n;++i) rank[i]=(unsigned char)s[i];
    for(int length=1;length<n;length<<=1) {
        auto key=[&](int position) {
            return pair<int,int>{
                rank[position],
                position+length<n?rank[position+length]:-1
            };
        };
        sort(order.begin(),order.end(),
             [&](int a,int b){ return key(a)<key(b); });
        new_rank[order[0]]=0;
        for(int i=1;i<n;++i)
            new_rank[order[i]]=new_rank[order[i-1]]
                +(key(order[i-1])!=key(order[i]));
        rank.swap(new_rank);
        if(rank[order.back()]==n-1) break;
    }
    return order;
}

vector<int> lcp_array(const string &s,const vector<int> &sa) {
    int n=s.size();
    assert((int)sa.size()==n);
    if(n<=1) return {};
    vector<int> rank(n),lcp(n-1);
    for(int i=0;i<n;++i) rank[sa[i]]=i;
    for(int position=0,length=0;position<n;++position) {
        int index=rank[position];
        if(index==n-1) {
            length=0;
            continue;
        }
        int next=sa[index+1];
        while(position+length<n&&next+length<n&&
              s[position+length]==s[next+length]) ++length;
        lcp[index]=length;
        if(length) --length;
    }
    return lcp;
}
\end{lstlisting}

Ejemplo de uso:

\begin{lstlisting}[style=cpp]
string s="banana";
vector<int> sa=suffix_array(s);  // 5 3 1 0 4 2
vector<int> lcp=lcp_array(s,sa); // 1 3 0 0 2
ll distinct=(ll)s.size()*(s.size()+1)/2
            -accumulate(lcp.begin(),lcp.end(),0LL);
cout<<distinct<<'\n'; // 15
// Cantidad de substrings distintos:
ll total=(ll)n*(n+1)/2;
ll repeated=accumulate(lcp.begin(),lcp.end(),0LL);
ll distinct=total-repeated;
//Porque hay \(n(n+1)/2\) substrings en total, pero cada lcp[i] representa prefijos ya vistos en el sufijo anterior.
\end{lstlisting}

La cantidad de substrings distintos es
\[
\frac{n(n+1)}2-\sum_i lcp[i].
\]
Para consultar repetidamente el LCP entre dos sufijos cualesquiera, guardar la
posición de cada sufijo en \texttt{sa} y construir un RMQ sobre el arreglo
\texttt{lcp}.